\documentclass[11pt]{article}

\usepackage[margin=1in]{geometry}
\usepackage{amsmath,amssymb}
\usepackage[hidelinks]{hyperref}

\newcommand{\R}{\mathbb{R}}
\newcommand{\otpi}{\widehat{\otimes}_{\pi}}
\newcommand{\oteps}{\widehat{\otimes}_{\epsilon}}
\newcommand{\FL}{\mathcal{F}L}
\newcommand{\Mod}{\mathcal{M}}

\title{A Literature-Implied Partial Answer to Problem 3.19}
\author{Run \texttt{fa\_banach\_001}, lane 1}
\date{2026-07-19}

\begin{document}
\maketitle

\section*{Status}

\textbf{Literature-implied answer, partial subcase.}
This note records an agent-identified implication of existing theorems. It
should not be counted as an original proof. The supporting paper does not
explicitly state that it answers Problem 3.19 of \cite{DP2023}; rather, its
tensor-product modulation-space identifications combine with the Gabor
coefficient theorem in \cite{DP2023} to give the sequence-space descriptions
below.

\section*{Source Problem}

The source paper is Debrouwere--Prangoski, \emph{Gabor frame characterisations
of generalised modulation spaces}, arXiv:2102.03217, published in
\emph{Analysis and Applications} \cite{DP2023}. In Section 3, Problem 3.19
asks:
\[
 \text{describe explicitly }
 (L^{p_1}\widehat{\otimes}_{\tau}L^{p_2})^{B_0}_d(\Lambda_1\times\Lambda_2),
 \qquad \tau\in\{\pi,\epsilon\},
\]
for lattices \(\Lambda_1,\Lambda_2\subset\R^n\) and
\(1\leq p_1,p_2<\infty\).

\section*{Identified Subcase}

Combining \cite[Corollary 4.11]{DP2023} with
\cite[Corollary 6.1 and Remark 6.2]{FPP2022} gives the following partial
description of Problem 3.19.

\begin{itemize}
\item If \(1\leq p_1\leq p_2\leq 2\), then
\[
 (L^{p_1}\otpi L^{p_2})^{B_0}_d(\Lambda_1\times\Lambda_2)
 \cong
 \ell^1(\Lambda_1;\ell^{p_2}(\Lambda_2))
\]
topologically.
\item If \(2\leq p_2\leq p_1<\infty\), then
\[
 (L^{p_1}\oteps L^{p_2})^{B_0}_d(\Lambda_1\times\Lambda_2)
 \cong
 c_0(\Lambda_1;\ell^{p_2}(\Lambda_2))
\]
topologically.
\end{itemize}

These formulae extend the special \(p_2=2\) descriptions stated in
\cite[Corollary 3.21]{DP2023}; the projective formula also contains the
source paper's explicitly mentioned trivial case \(p_1=1\).

\section*{Reason For The Identification}

Let \(F=L^{p_1}\otpi L^{p_2}\). The supporting tensor-product paper
\cite[Corollary 6.1(a) and Remark 6.2]{FPP2022} proves
\[
 \Mod[F] = W(\FL^{p_2},L^1)=\mathcal{F}M^{p_2,1},
 \qquad 1\leq p_1\leq p_2\leq 2.
\]
The source paper's Gabor coefficient theorem
\cite[Corollary 4.11]{DP2023} identifies \(\Mod[F]\), through the coefficient
map associated with a suitable dual Gabor pair, with \(F_d^{\sigma}\), where
\(F_d^{\sigma}=F_d^{B_0}\). The standard Gabor coefficient description of
\(\mathcal{F}M^{p_2,1}=W(\FL^{p_2},L^1)\) is
\(\ell^1(\Lambda_1;\ell^{p_2}(\Lambda_2))\). This yields the projective
formula.

Similarly, for \(F=L^{p_1}\oteps L^{p_2}\),
\cite[Corollary 6.1(b)]{FPP2022} proves
\[
 \Mod[F] = W(\FL^{p_2},L^\infty_0),
 \qquad 2\leq p_2\leq p_1<\infty.
\]
Applying \cite[Corollary 4.11]{DP2023} and the standard Gabor coefficient
description of \(W(\FL^{p_2},L^\infty_0)\) gives
\(c_0(\Lambda_1;\ell^{p_2}(\Lambda_2))\).

For arbitrary product lattices, one may pass to a sufficiently fine finite-index
product refinement that admits the required smooth dual Gabor pair, apply the
preceding identification there, and restrict by zero extension to the original
sub-lattice. Since the definition of \(E_d^B\) is independent of the compact
test function used in \cite[Theorem 3.1]{DP2023}, this restriction preserves
the corresponding mixed sequence-space description.

\section*{Scope Limitations}

This does not settle Problem 3.19 for all exponents. The remaining regions
include the projective tensor product outside \(1\leq p_1\leq p_2\leq 2\) and
the injective tensor product outside \(2\leq p_2\leq p_1<\infty\). The
supporting authors did not explicitly advertise their result as an answer to
Problem 3.19; the relation is a reformulation through the source paper's Gabor
coefficient theorem.

\section*{Verification Notes}

The cheap run indexes were searched for arXiv:2102.03217, the title, Gabor
frames, generalized modulation spaces, and tensor-product \(L^p\) discrete
spaces. No previous packet for this source paper was found. Web and local
source searches found the published source metadata and the related
Feichtinger--Pilipovic--Prangoski tensor-product modulation paper, but no
separate explicit statement of the above \(B_0\)-discrete-space answer.

\begin{thebibliography}{9}
\bibitem{DP2023}
A. Debrouwere and B. Prangoski,
\emph{Gabor frame characterisations of generalised modulation spaces},
Analysis and Applications 21 (2023), no. 3, 547--596.
arXiv:2102.03217. DOI:
\href{https://doi.org/10.1142/S0219530522500178}{10.1142/S0219530522500178}.

\bibitem{FPP2022}
H. G. Feichtinger, S. Pilipovic, and B. Prangoski,
\emph{Modulation spaces associated with tensor products of amalgam spaces},
Annali di Matematica Pura ed Applicata 201 (2022), 127--155.
arXiv:2012.12295. DOI:
\href{https://doi.org/10.1007/s10231-021-01110-9}{10.1007/s10231-021-01110-9}.
\end{thebibliography}

\end{document}
